\section{Discussion} \label{sec:discussion}
The original premise of this work is that probe-to-gallery retrieval in heterogeneous camera networks is direction-dependent: a clean query retrieved against a degraded gallery is not the same problem as the reverse. The Randers drift factor $\hat{\mathbf{z}}^{\omega}$ was intended to encode this per-camera directional shift: larger domain shifts should correspond to larger drift magnitudes, enabling the asymmetric scoring term eq.~\eqref{eq:finsler_distance} to compensate for camera-specific degradation.

The best-performing recipe contradicts this premise directly. Cross-camera drift coherence regularization ($\mathcal{L}_{\mathrm{dcc}}$) penalizes \emph{dispersion} of drift vectors across cameras for the same identity: it forces drift to be \emph{camera-invariant} for a given person. This is the opposite of what the motivating intuition would prescribe. Two complementary interpretations are consistent with the empirical data:

\textbf{(i) Drift as an identity-conditioned directional bias.} $\mathcal{L}_{\mathrm{dcc}}$ does not suppress drift magnitude; post-GS norms range $0.054$--$0.067$ across corruption severities and are monotonically increasing with severity ($\rho{=}0.67$, measured on the Arm~1b drift magnitude $\|\hat{\mathbf{z}}^{\omega}\|_2$; the toy study's severity-encoding scalar $\theta$ achieves $\rho\!\approx\!-0.91$ on the toy split under a different architecture and training protocol, see Sec.~\ref{subsec:controlled_corruption}). What $\mathcal{L}_{\mathrm{dcc}}$ suppresses is \emph{cross-camera variance}: drift for a given identity converges to a stable direction regardless of camera. Under this interpretation, $\hat{\mathbf{z}}^{\omega}$ encodes a per-person directional preference in feature space, indicating a stable asymmetric bias, not a camera noise sink. The empirical observation that identity distance grows monotonically with severity ($0 {\to} 1.27$) while drift grows only weakly ($0.054{\to}0.067$) is consistent: corruption bleeds into the identity slice, while drift remains bounded and coherent.

\textbf{(ii) Drift coherence as training-geometry regularization.} Bidirectional Finsler triplet mining requires that both $d_F(i,j)$ and $d_F(j,i)$ provide a reliable margin signal. When drift vectors are incoherent across cameras, the asymmetric term introduces high-variance distance estimates that destabilize hard-pair mining. $\mathcal{L}_{\mathrm{dcc}}$ functions as a variance regularizer that makes bidirectional mining numerically stable, explaining why the bidirectional gain ($+1.0\,\mathrm{pp}$ mAP) materializes only when $\mathcal{L}_{\mathrm{dcc}}$ is active (Arm~1b vs.\ Arms~0, 3, 4 in Table~\ref{tab:aux_loss_sweep}).

Both interpretations are consistent with $\Delta\mathrm{mAP}{\approx}0$ when switching from Euclidean to Finsler ranking on the same checkpoint (Supp.~\ref{sec:suppl_eval_drift_tables}): the gain from the primary recipe comes through \emph{training geometry}, not a fundamentally different test-time scoring function.

The toy study (Sec.~\ref{subsec:controlled_corruption}) mechanistically complements the main training stack. The full Finsler BAU system leaves the origin of its marginal gain ambiguous (training geometry or test-time Randers scoring; see Supp.~\ref{sec:suppl_eval_drift_tables}), whereas the toy setting separates the two: $\mathcal{L}_{\mathrm{mono}}$ reliably trains a severity-monotone scalar (Sec.~\ref{subsec:controlled_corruption}), the disabled-loss control causally isolates the mechanism, and the Randers gap is confirmed geometrically on balanced pairs. In the full system, $\mathcal{L}_{\mathrm{dcc}}$ plays the analogous role: both losses structure the drift subspace during training; $\mathcal{L}_{\mathrm{mono}}$ imposes severity monotonicity on a scalar projection while $\mathcal{L}_{\mathrm{dcc}}$ bounds cross-camera variance on a high-dimensional drift vector. Neither guarantees that test-time Randers scoring outperforms Euclidean ranking on the same checkpoint, but $\mathcal{L}_{\mathrm{mono}}$ additionally enables the directional gap by construction.
